\providecommand{\Needspace}[1]{}
\Needspace{8\baselineskip}
\item \textbf{Clasificación de mediciones de brillo con funciones y bucles}.
\nopagebreak[4]

\begin{tcolorbox}[colback=blue!2!white,colframe=blue!55!black,title={Enunciado},
                  before skip=3pt,after skip=5pt,boxsep=3pt,top=3pt,bottom=3pt,
                  breakable]
\small
\raggedright
Durante una noche de observación se registran valores de flujo relativo de una
estrella. Para resumir las mediciones se usará una función, un ciclo y
condicionales.

Use la siguiente lista:
\[
\texttt{flujos = [0.82,\ 1.05,\ 1.18,\ 0.76,\ 1.34,\ 0.91,\ 1.27,\ 0.69]}.
\]

Una medición se clasifica según su flujo \(F\):
\[
F < 0{.}85 \Longrightarrow \text{baja},
\qquad
0{.}85 \leq F \leq 1{.}15 \Longrightarrow \text{normal},
\qquad
F > 1{.}15 \Longrightarrow \text{alta}.
\]

Realice las siguientes operaciones:

\begin{itemize}
\setlength{\topsep}{4pt}
\setlength{\itemsep}{3pt}
\setlength{\parsep}{0pt}
\item[(a)] Defina la lista \texttt{flujos}.
\item[(b)] Cree una función llamada \texttt{clasificar\_flujo(F)} que reciba un
flujo y retorne el texto \texttt{"baja"}, \texttt{"normal"} o \texttt{"alta"}
según los criterios anteriores.
\item[(c)] Recorra la lista con un ciclo \texttt{for}. Para cada flujo, use la
función creada, imprima el índice, el flujo y su clasificación.
\item[(d)] Cuente cuántas mediciones quedan en cada categoría. Use contadores
simples que aumenten dentro del ciclo.
\item[(e)] Calcule el promedio de los flujos usando una suma acumulada y
\texttt{len(flujos)}.
\item[(f)] Use un condicional final para indicar si la noche fue estable o no:
si la cantidad de mediciones normales es mayor o igual que la suma de mediciones
bajas y altas, imprima \texttt{"noche estable"}; en caso contrario, imprima
\texttt{"noche variable"}.
\end{itemize}
\end{tcolorbox}

\begin{solution}
\begin{pythonjupyter}
flujos = [0.82, 1.05, 1.18, 0.76, 1.34, 0.91, 1.27, 0.69]

def clasificar_flujo(F):
    if F < 0.85:
        return "baja"
    elif F <= 1.15:
        return "normal"
    else:
        return "alta"

cantidad_baja = 0
cantidad_normal = 0
cantidad_alta = 0
suma_flujos = 0

for i in range(len(flujos)):
    F = flujos[i]
    clasificacion = clasificar_flujo(F)
    suma_flujos = suma_flujos + F

    if clasificacion == "baja":
        cantidad_baja = cantidad_baja + 1
    elif clasificacion == "normal":
        cantidad_normal = cantidad_normal + 1
    else:
        cantidad_alta = cantidad_alta + 1

    print("Índice:", i)
    print("Flujo:", F)
    print("Clasificación:", clasificacion)
    print("-" * 20)

promedio = suma_flujos / len(flujos)

print("Cantidad baja:", cantidad_baja)
print("Cantidad normal:", cantidad_normal)
print("Cantidad alta:", cantidad_alta)
print("Promedio:", promedio)

if cantidad_normal >= cantidad_baja + cantidad_alta:
    print("noche estable")
else:
    print("noche variable")
\end{pythonjupyter}
\end{solution}
